%% ============================================================
%%  Crime, Law and Social Change  |  Springer  |  ISSN 0925-4994
%%  Special Issue: AI Enabled Security Frameworks for
%%                 Emerging Cybercrime and Digital Threats
%%
%%  COMPILE ON OVERLEAF: article class — compiles without any
%%  extra files.  For final Springer submission, switch to
%%  \documentclass[runningheads]{svjour3} and upload the
%%  svjour3.cls pack from springer.com/authors.
%%
%%  Compile order: pdflatex -> bibtex -> pdflatex -> pdflatex
%% ============================================================
\documentclass[12pt,a4paper]{article}

%% ── Page geometry ────────────────────────────────────────────
\usepackage[top=2.5cm,bottom=2.5cm,left=3cm,right=3cm]{geometry}

%% ── Fonts and encoding ───────────────────────────────────────
\usepackage[T1]{fontenc}
\usepackage[utf8]{inputenc}
\usepackage{lmodern}
\usepackage{microtype}

%% ── Mathematics ──────────────────────────────────────────────
\usepackage{amsmath,amssymb,amsthm}

%% ── Graphics ─────────────────────────────────────────────────
\usepackage{graphicx}
\usepackage{subcaption}
\graphicspath{{figures/}}

%% ── Tables ───────────────────────────────────────────────────
\usepackage{booktabs}
\usepackage{multirow}
\usepackage{tabularx}
\usepackage{array}

%% ── Algorithms ───────────────────────────────────────────────
\usepackage[ruled,vlined,linesnumbered]{algorithm2e}

%% ── Section formatting ───────────────────────────────────────
\usepackage{titlesec}
\titleformat{\section}{\normalfont\large\bfseries}{\thesection}{1em}{}
\titleformat{\subsection}{\normalfont\normalsize\bfseries}{\thesubsection}{1em}{}
\titleformat{\subsubsection}{\normalfont\normalsize\itshape}{\thesubsubsection}{1em}{}

%% ── Header / footer ──────────────────────────────────────────
\usepackage{fancyhdr}
\pagestyle{fancy}
\fancyhf{}
\fancyhead[L]{\small\textit{Crime, Law and Social Change}}
\fancyhead[R]{\small\textit{Korde: MixTrace CoinJoin Deanonymization}}
\fancyfoot[C]{\thepage}
\renewcommand{\headrulewidth}{0.4pt}

%% ── Bibliography ─────────────────────────────────────────────
\usepackage[round,authoryear]{natbib}
\bibliographystyle{spbasic}

%% ── Hyperlinks ───────────────────────────────────────────────
\usepackage[colorlinks=true,linkcolor=blue,
            citecolor=blue,urlcolor=blue,
            pdftitle={Solving the CoinJoin Mystery},
            pdfauthor={Sagar Korde}]{hyperref}
\usepackage{url}

%% ── Line spacing ─────────────────────────────────────────────
\usepackage{setspace}
\onehalfspacing

%% ── Abstract formatting ──────────────────────────────────────
\usepackage{abstract}
\renewcommand{\abstractnamefont}{\normalfont\bfseries}
\renewcommand{\abstracttextfont}{\normalfont\small}

%% ── Custom commands ──────────────────────────────────────────
\DeclareMathOperator*{\argmin}{arg\,min}
\DeclareMathOperator{\Entropy}{H}
\newcommand{\eg}{for instance}
\newcommand{\ie}{that is}

%% ── Title metadata ───────────────────────────────────────────
\title{\textbf{Solving the CoinJoin Mystery: A Multi-Modal Graph
       Attention Framework for Bitcoin Transaction
       Deanonymization}%
       \thanks{The author declares no competing interests.
               No external funding was received.}}

\author{Sagar Korde\\
  \small Independent Researcher, Pune, India\\
  \small \href{mailto:sagarkorde04@gmail.com}{sagarkorde04@gmail.com}}

\date{\small Submitted: \today}

% ─────────────────────────────────────────────────────────────
\begin{document}
% ─────────────────────────────────────────────────────────────
\maketitle
\thispagestyle{fancy}

% ─────────────────────────────────────────────────────────────
%  ABSTRACT
% ─────────────────────────────────────────────────────────────
\begin{abstract}
\noindent
Bitcoin CoinJoin transactions deliberately obscure payment flows
by merging unrelated users' unspent transaction outputs into a
single on-chain record, substantially weakening the
Common Input Ownership Heuristic that underpins most operational
blockchain forensic tools. Existing detection pipelines suffer
from three compounding limitations: structural address clustering
fails under Taproot's Pay-to-Taproot script indistinguishability,
CoinJoin Spending Transaction linking relies on one-dimensional
temporal proximity while ignoring structural signals, and
supervised graph classifiers produce poor minority-class recall
on the severely class-imbalanced Elliptic Bitcoin Dataset.
The present study introduces MixTrace, a multi-modal framework
that addresses all three limitations in a unified experimental
pipeline. A Taproot-Aware DBSCAN clustering strategy reduces
the false-positive rate of naive heuristic classification by
79.8 percent across 200,000 analysed transactions, achieving a
silhouette score of 0.973 across the derived structural clusters.
A four-dimensional Chamfer Distance formulation extends temporal
CoinJoin Spending Transaction linking to the joint space of
timestamp, fee rate, input count, and output count. Deep Graph
Infomax pre-training on all 203,769 Elliptic graph nodes,
followed by Graph Attention Network fine-tuning with class-weighted
cross-entropy, achieves an illicit-class F1-score of 0.7337,
precision of 0.7094, recall of 0.7598, and ROC-AUC of 0.9868
on the chronological held-out test set, outperforming Random
Forest (F1: 0.6285), K-Nearest Neighbours (F1: 0.4829),
K-Means (F1: 0.0883), and a standalone Graph Attention Network
without pre-training (F1: 0.2037). An ablation study confirms
that the Deep Graph Infomax skip connection contributes 36
percentage points of the final F1-score, validating the use
of self-supervised pre-training on unlabelled blockchain nodes.
The complete reproducible 14-step pipeline is publicly available
to support forensic and compliance research communities.

\smallskip\noindent\textbf{Keywords:} Bitcoin; CoinJoin;
Graph Attention Network; Deep Graph Infomax; Blockchain forensics;
Chamfer distance; Illicit transaction detection; ADASYN oversampling
\end{abstract}

\newpage

% ─────────────────────────────────────────────────────────────
\section{Introduction}
\label{sec:intro}
% ─────────────────────────────────────────────────────────────

The rapid globalisation of cryptocurrency-facilitated crime
has placed blockchain forensics at the centre of contemporary
financial intelligence practice. Ransomware operators, darknet
market participants, and sanctions evaders collectively laundered
an estimated \$23.8 billion in cryptocurrency during 2022 alone,
with Bitcoin remaining the dominant settlement medium for illicit
flows~\citep{lo2023inspection,ahmed2024federated}. Bitcoin's
public ledger records every transaction in a pseudonymous but
auditable form, creating a structural tension between user
privacy and investigative transparency that has driven a decade
of forensic research~\citep{motie2024gnn}.

The primary clustering tool in blockchain forensics is the
Common Input Ownership Heuristic (CIOH), which assumes that
all input addresses co-signing a transaction belong to a single
controlling entity. CoinJoin~\citep{ouyang2024bitcoin} disrupts
this assumption by enabling multiple independent parties to
collaboratively construct a single transaction that merges their
unspent transaction outputs (UTXOs), breaking the CIOH and
severing the fund-flow trail that investigators rely upon.
Privacy-focused wallets such as Wasabi and JoinMarket have
processed hundreds of millions of dollars through CoinJoin
coordination, necessitating purpose-built deanonymization
approaches beyond simple heuristic clustering~\citep{coinjoin2024ieee}.

Three research gaps persist in the current literature that
collectively limit the practical effectiveness of automated
CoinJoin analysis.

\textbf{Gap 1 -- Taproot Clustering Failure.} Bitcoin's Taproot
soft fork, activated at block height 709,632 in November 2021,
introduced Pay-to-Taproot (P2TR) scripts based on Schnorr
signatures and key aggregation through MuSig2 protocol. The
critical forensic consequence is that a single-user P2TR output
is computationally indistinguishable on-chain from a multi-party
aggregated P2TR output used in CoinJoin coordination. Naive CIOH
clustering applied to Taproot transactions therefore generates
systematic false positives by incorrectly flagging legitimate
single-user P2TR activity as mixing events. Studies of
post-Taproot address clustering have consistently identified this
as an emerging forensic blind spot~\citep{lee2024censor,asiri2025gcn}.

\textbf{Gap 2 -- One-Dimensional CoinJoin Spending Transaction Linking.}
Deanonymizing post-mix spending requires linking the outputs of
a mixing event to the subsequent spending transactions of
individual participants. The only published similarity metric
for this task applies a one-dimensional Chamfer distance
over transaction timestamps~\citep{chen2025multidist}. Fee-rate
behaviour, input multiplicity, and output cardinality carry
independent structural signals that remain unexploited in
current linking pipelines, leaving measurable discriminative
information on the table.

\textbf{Gap 3 -- Class Imbalance in Graph-Based Illicit Detection.}
The Elliptic Bitcoin Dataset~\citep{fan2024gnn}, the canonical
benchmark for illicit node detection, contains 4,545 illicit
and 42,019 licit labelled transactions, yielding a 9.2:1
licit-to-illicit imbalance ratio. Standard graph classifiers
suppress minority-class predictions to maximise overall accuracy,
producing F1-scores below 0.50 for the illicit class in most
published configurations~\citep{alarab2024recurrent}. The 157,205
unlabelled nodes in the same graph represent a large reservoir
of structural information that supervised methods discard.

To address these three gaps in a unified experimental framework,
the present study introduces \textbf{MixTrace}, comprising three
coordinated components: (i) a Taproot-Aware DBSCAN clustering
strategy on multi-dimensional structural and behavioural transaction
features; (ii) a Multi-Dimensional Chamfer Distance (MD-CD)
formulation that extends CST linking from one dimension to
four dimensions; and (iii) a two-phase graph learning architecture
applying Deep Graph Infomax (DGI) self-supervised pre-training
before supervised Graph Attention Network fine-tuning with
class-weighted cross-entropy and ADASYN oversampling.

Experiments are conducted on the public Elliptic Bitcoin Dataset
and an author-curated corpus of 5,884,387 Bitcoin transactions
enriched with nine novel engineered features. Results demonstrate
that MixTrace achieves an illicit-class F1-score of 0.7337 and
ROC-AUC of 0.9868, representing improvements of 10.5 percentage
points in F1 over the best conventional baseline and 79.8 percent
reduction in clustering false-positive rate.

The manuscript is structured as follows. Section~\ref{sec:related}
reviews the relevant literature. Section~\ref{sec:method} presents
the MixTrace methodology. Section~\ref{sec:exp} describes
experimental setup and results. Section~\ref{sec:discuss} analyses
findings. Section~\ref{sec:conclude} summarises contributions
and directions for future research.

% ─────────────────────────────────────────────────────────────
\section{Related Work}
\label{sec:related}
% ─────────────────────────────────────────────────────────────

\subsection{Graph Neural Networks for Blockchain Forensics}

Graph neural network (GNN) architectures have emerged as the
dominant paradigm for illicit transaction detection on public
blockchains. \citet{asiri2025gcn} applied graph convolutional
networks to Bitcoin transaction graphs and demonstrated that
multi-hop neighbourhood aggregation captures structural patterns
invisible to transaction-level feature classifiers.
\citet{alarab2024recurrent} proposed a robust recurrent GCN
architecture that integrates sequential temporal modelling with
spatial message passing, achieving improved stability on the
Elliptic benchmark compared to static GCN baselines.
\citet{ouyang2024bitcoin} framed Bitcoin anti-money laundering
as a subgraph contrastive learning problem, constructing
positive and negative subgraph pairs to train a representation
that clusters illicit transaction neighbourhoods distinctly
from licit flows.
\citet{chen2025multidist} introduced a multi-distance
spatial-temporal GNN for blockchain anomaly detection that
exploits transaction timing patterns at multiple temporal
resolutions, complementing the single-resolution temporal
feature used in the present study's Chamfer distance baseline.
The global-local graph attention framework of~\citet{globalloc2025}
combines local multi-head attention with a global pooling
context vector, using cyclic pseudo-labelling to propagate
supervisory signal into the unlabelled node majority.

\subsection{Anti-Money Laundering with Machine Learning}

The relationship between graph representation learning and
financial fraud detection has been thoroughly surveyed by
\citet{motie2024gnn}, who conducted a systematic review of 94
primary studies and identified class imbalance, temporal
graph evolution, and the exploitation of unlabelled data as
the three most critical open challenges. \citet{fan2024gnn}
similarly reviewed GNN methods for financial fraud, concluding
that attention-based neighbourhood aggregation consistently
outperforms isotropic mean aggregation when fraudulent and
benign subgraph patterns differ structurally.
\citet{xiao2025biEGCN} proposed Balanced-BiEGCN, a bidirectional
evolving GCN with class-balanced learning for dynamic Bitcoin
anomaly detection that addresses temporal label drift, an
issue particularly relevant to the Elliptic dataset's
49-time-step span.
\citet{lee2024censor} introduced CENSor, a GCN-based hyperedge
classification framework that groups Bitcoin transactions into
operation-level hyperedges before assigning illicit scores,
achieving precision improvements over node-level classifiers.
\citet{lo2023inspection} conducted a benchmark comparison of
machine learning methods for blockchain anti-money laundering,
identifying that graph-structural features derived from the
directed edge list provide the most consistent performance
improvement across evaluated classifiers.

\subsection{CoinJoin Privacy and Deanonymization}

The forensic analysis of CoinJoin transactions occupies a
specialised sub-domain of blockchain privacy research.
\citet{coinjoin2024ieee} studied the quantitative impact of
CoinJoin participation on Bitcoin network traceability,
measuring reduction in CIOH clustering confidence as a
function of mixing round count and wallet-level behaviour.
Anomaly detection methods applied to cryptocurrency networks
have been reviewed by~\citet{futureinternet2025}, who identified
DBSCAN and isolation-forest variants as the most effective
unsupervised baselines for exchange-level flagging.
\citet{kiani2024smartcontract} examined the structural
properties of non-standard Bitcoin scripts including multi-signature
coordination patterns that overlap with CoinJoin outputs,
noting that post-Taproot script type convergence creates
classification challenges for rule-based detection systems.
Phishing detection research on Ethereum by~\citet{chen2024sensors}
demonstrated that data augmentation combined with hybrid GNN
architectures substantially reduces false positives caused by
structural overlap between legitimate and malicious transaction
patterns, a finding directly applicable to the Taproot
indistinguishability problem addressed in the present study.
\citet{choi2024electronics} showed that transformer-based
traversal of transaction graphs improves detection of phishing
scams by capturing long-range transaction path dependencies
that local message-passing GNNs miss.

\subsection{Self-Supervised and Contrastive Graph Learning}

Deep Graph Infomax (DGI)~\citep{velickovic2019deep} frames
unsupervised graph representation learning as a mutual
information maximisation problem between local node embeddings
and a global graph summary, enabling pre-training on unlabelled
graph data.
\citet{graphanomaly2024} applied graph anomaly detection with
disentangled prototypical autoencoders to cryptocurrency
transaction graphs, demonstrating that self-supervised
prototype learning on unlabelled transactions effectively
captures the latent structure of illicit behavioural clusters.
\citet{wei2025temporal} proposed a temporal-aware heterogeneous
graph oversampling strategy for Internet fraud detection that
combines attention-based neighbourhood fusion with synthetic
minority sample generation, addressing both class imbalance
and temporal heterogeneity simultaneously.
Federated learning extensions of blockchain graph classifiers
have been investigated by~\citet{ahmed2024federated}, who
demonstrated that privacy-preserving collaborative training
across multiple exchange nodes maintains detection accuracy
competitive with centralised graph models.
\citet{baabdullah2024federated} further studied the interaction
between federated training and class-balanced loss functions
in credit card fraud detection systems, finding that federated
aggregation preserves class-weighting effectiveness when the
participating clients share similar imbalance ratios.

\subsection{Class Imbalance Mitigation in Fraud Detection}

Handling severe label imbalance is a persistent challenge in
financial fraud detection. \citet{tekkali2024adasyn} proposed
an advancement of the ADASYN algorithm with density-sensitive
synthesis boundaries, demonstrating improved F1-score on
fraud detection benchmarks with imbalance ratios up to 100:1.
\citet{creditcard2024adasyn} evaluated ADASYN combined with
SHAP-based feature importance for credit card fraud detection,
finding that selective feature retention before oversampling
reduces synthetic sample noise and improves classifier
calibration.
\citet{albalawi2025credit} compared traditional and deep
learning models under class imbalance mitigation strategies
including ADASYN, SMOTE, and class weighting, concluding
that inverse-frequency class weighting within the loss
function produces the most stable precision-recall trade-off
for neural network architectures because it preserves the
original sample geometry rather than introducing synthetic
points.
\citet{chithanuru2025ethereum} and \citet{gu2025ensemble}
respectively validated XAI-enabled ensemble methods and
stacking ensembles for Ethereum fraud detection, both noting
that the combined use of minority-class oversampling during
training and threshold optimisation at inference substantially
improves F1-score on imbalanced blockchain benchmarks.
The present study adopts inverse-frequency class weighting as
the primary imbalance strategy in the graph loss function,
with ADASYN applied separately as a diagnostic reference
consistent with the recommendations of~\citet{tekkali2024adasyn}.

\subsection{Explainability and Benchmarking}

Interpretability of blockchain fraud models has gained
increasing regulatory relevance. \citet{ethereum2025xai}
developed a near-real-time Ethereum fraud detection framework
with integrated explainability, finding that gradient-based
attribution methods identify transaction fee and timing
features as primary drivers of illicit classification in GCN
models, consistent with the novel temporal irregularity feature
engineered in the present study.
The heterogeneous graph transformer for Ethereum phishing
detection introduced by~\citet{wei2025temporal} established
that multi-relational graph encoding captures complementary
signals from address-to-transaction and transaction-to-address
edge types, motivating the bidirectional edge construction
used in MixTrace's PyTorch Geometric graph object.

% ─────────────────────────────────────────────────────────────
\section{Methodology}
\label{sec:method}
% ─────────────────────────────────────────────────────────────

\subsection{Problem Formulation}

Let $\mathcal{G} = (\mathcal{V}, \mathcal{E}, \mathbf{X})$
denote the Elliptic transaction graph, where $\mathcal{V}$
is the set of $N = 203{,}769$ transaction nodes, $\mathcal{E}$
the set of 234,355 directed edges representing Bitcoin payment
flows, and $\mathbf{X} \in \mathbb{R}^{N \times F}$ the
node feature matrix with $F = 169$ dimensions (165 original
Elliptic features augmented with four novel graph-structural
features introduced in Section~\ref{sec:features}).
Each node $v \in \mathcal{V}$ carries a discrete time step
$\tau(v) \in \{1,\ldots,49\}$ and a label
$y_v \in \{1, 0, -1\}$, where 1 denotes illicit,
0 denotes licit, and $-1$ denotes unknown.
The labelled subset comprises 46,564 nodes; the remaining
157,205 carry no label.
The learning objective is binary node classification on the
labelled subset, evaluated on nodes with
$\tau(v) \in \{42,\ldots,49\}$ using precision, recall,
F1-score, PR-AUC, and ROC-AUC on the illicit class.

\subsection{Novel Feature Engineering}
\label{sec:features}

Nine novel features are computed across the two experimental
datasets to supplement existing representations.

\subsubsection{Author Dataset: Five Transaction-Level Features}

\paragraph{Denomination Entropy.}
CoinJoin transactions characteristically produce multiple
outputs of equal value to obscure the payment-to-change
distinction. For transaction $t$ with output value vector
$\mathbf{v} = (v_1,\ldots,v_k)$, the normalised output
probability vector is $p_i = v_i / \sum_j v_j$, and the
denomination entropy is:
\begin{equation}
  \Entropy(t) = -\sum_{i=1}^{k} p_i \log_2 p_i.
  \label{eq:entropy}
\end{equation}
High entropy ($\Entropy \to \log_2 k$) signals equal-valued
outputs, a structural CoinJoin signature.

\paragraph{Mixing Index.}
The combinatorial ambiguity of input-output pairing is
measured by:
\begin{equation}
  \text{MI}(t) = \frac{|\mathcal{A}| \cdot |\mathcal{B}|}{|\mathcal{A} \cup \mathcal{B}|},
\end{equation}
where $\mathcal{A}$ and $\mathcal{B}$ are the input and output
address sets. In the author dataset, MI has a mean of 8.79
and standard deviation of 689.27, reflecting the presence of
high-fan-out mixing transactions at the upper tail.

\paragraph{Temporal Irregularity.}
Mixing coordination services schedule transactions at atypical
network hours to reduce mempool competition.
Temporal irregularity is the $z$-score of a transaction's
fee rate within the hour-of-day and day-of-week cell:
\begin{equation}
  \text{TI}(t) = \frac{f_t - \mu_{h,d}}{\sigma_{h,d}},
\end{equation}
where $f_t$ is the fee rate in sat/vbyte and $\mu_{h,d}$,
$\sigma_{h,d}$ are the cell mean and standard deviation.

\paragraph{UTXO Concentration.}
A Herfindahl-Hirschman Index approximation over output values:
\begin{equation}
  \text{UC}(t) = r^2,
\end{equation}
where $r$ is the value-concentration ratio. CoinJoin outputs
are expected to produce UC near zero; single-payment
transactions yield UC near one.

\paragraph{CoinJoin Proximity Score.}
A weighted rule-based composite aggregating structural
CoinJoin indicators:
\begin{equation}
  \text{CPS}(t) = \sum_{r \in \mathcal{R}} w_r \cdot \mathbf{1}[t \models r],
\end{equation}
where $\mathcal{R}$ includes equal-denomination output flags,
input count exceeding three, zero address reuse, and round
output values. In the author dataset CPS has a mean of 0.2663
and standard deviation of 0.1307.

\subsubsection{Elliptic Dataset: Four Graph-Structural Features}

\paragraph{Graph In-Degree and Out-Degree.}
In-degree and out-degree are computed directly from the
directed edge list. Mean in-degree is 1.150 (standard
deviation 3.911) and mean out-degree is 1.150 (standard
deviation 1.895), with maximum values of 284 and 472
indicating a strongly skewed hub structure.

\paragraph{Illicit Neighbour Ratio.}
For each labelled node $v$, the fraction of labelled one-hop
neighbours carrying the illicit label:
\begin{equation}
  \text{INR}(v) = \frac{\sum_{u \in \mathcal{N}(v)} \mathbf{1}[y_u = 1]}
                       {|\{u \in \mathcal{N}(v) : y_u \neq -1\}|}.
\end{equation}
INR achieves a Pearson correlation of $r = +0.4159$ with
the ground-truth label.

\paragraph{Shortest Path to Illicit.}
A breadth-first search from each confirmed illicit source
node computes the minimum hop distance to every reachable
node, capped at six. This feature achieves a Pearson
correlation of $r = -0.9182$ with the illicit label,
the strongest single-feature predictor in the augmented
feature set, consistent with the guilt-by-association
principle documented in financial fraud network analysis
by~\citet{gu2025ensemble}.

\subsection{Multi-Dimensional Chamfer Distance}
\label{sec:mdcd}

The one-sided Chamfer distance from point set
$\mathcal{U} = \{u_1,\ldots,u_m\}$ to
$\mathcal{V} = \{v_1,\ldots,v_n\}$ in $\mathbb{R}^d$ is:
\begin{equation}
  d_C(\mathcal{U},\mathcal{V})
    = \frac{1}{|\mathcal{U}|} \sum_{u_i \in \mathcal{U}}
      \min_{v_j \in \mathcal{V}} \|u_i - v_j\|_2.
  \label{eq:chamfer}
\end{equation}
The baseline 1-D formulation uses only the Unix timestamp
($d=1$). The Multi-Dimensional Chamfer Distance sets $d=4$
with the standardised feature vector:
\begin{equation}
  \phi(t) = \bigl[
    \text{timestamp},\;
    \text{fee\_rate},\;
    \text{input\_count},\;
    \text{output\_count}
  \bigr]^\top,
  \label{eq:phi}
\end{equation}
where each dimension is independently normalised to zero
mean and unit variance prior to distance computation.
Transactions are grouped into one-hour temporal windows
indexed by block height divided by six. A candidate pair
$(W_i, W_j)$ is predicted as a linked CoinJoin-spend pair
if $d_C(W_i, W_j) < \theta^*$, where $\theta^*$ maximises
F1-score on a held-out validation partition.

\subsection{Taproot-Aware DBSCAN Clustering}
\label{sec:taproot_method}

\subsubsection{Naive CIOH Baseline}
The naive Common Input Ownership Heuristic flags every
multi-input transaction as a CoinJoin candidate:
\begin{equation}
  \hat{y}_\text{CIOH}(t) =
    \mathbf{1}[|\mathcal{A}| > 2] \cdot
    \mathbf{1}[|\mathcal{B}| > 2].
\end{equation}
This rule achieves perfect recall at the cost of systematic
false positives that increase with P2TR adoption.

\subsubsection{Taproot-Aware DBSCAN}
The proposed method clusters transactions using DBSCAN on
the six-dimensional structural feature vector:
\begin{equation}
  \psi(t) = \bigl[
    \text{input\_count},\;
    \text{output\_count},\;
    \text{fee\_rate},\;
    \text{value\_concentration\_ratio},\;
    \text{address\_reuse},\;
    \text{mixing\_index}
  \bigr]^\top.
  \label{eq:psi}
\end{equation}
All features are standardised before clustering.
DBSCAN is fitted with $\varepsilon = 0.6$ and
$\text{min\_samples} = 5$ using a $k$-d tree index
to maintain memory tractability~\citep{ester1996density}.
To address the $O(n^2)$ memory complexity at dataset scale,
DBSCAN is fitted on a stratified subsample of 30,000
transactions; cluster labels are propagated to the full
200,000-transaction analysis set via nearest-centroid
assignment within the $\varepsilon$ radius. Points falling
outside all cluster radii receive the naive CIOH prediction
as a fallback. Cluster CoinJoin labels are assigned by
majority vote over ground-truth labels of cluster members.

\subsection{MixTrace: DGI Pre-Training and GAT Fine-Tuning}
\label{sec:mixtrace_method}

\subsubsection{Phase 1: Deep Graph Infomax Pre-Training}

The graph encoder $f_\theta: \mathbb{R}^{N \times F}
\to \mathbb{R}^{N \times 128}$ is a two-layer GCN:
\begin{align}
  \mathbf{H}^{(1)} &= \text{ELU}\!\bigl(
    \hat{\mathbf{A}}\,\mathbf{X}\,\mathbf{W}_1\bigr)
    \in \mathbb{R}^{N \times 256}, \\
  \mathbf{Z} &= \text{ELU}\!\bigl(
    \hat{\mathbf{A}}\,\mathbf{H}^{(1)}\,\mathbf{W}_2\bigr)
    \in \mathbb{R}^{N \times 128},
\end{align}
where $\hat{\mathbf{A}}$ is the symmetrically normalised
adjacency matrix with self-loops. The global graph summary is
$\mathbf{s} = \sigma(\text{mean}_v\, \mathbf{z}_v) \in
\mathbb{R}^{128}$. A bilinear discriminator
$\mathcal{D}(\mathbf{z}_i, \mathbf{s}) =
\mathbf{z}_i^\top \mathbf{W}_d\, \mathbf{s}$
scores the agreement between a node embedding and the summary.
Corrupted embeddings $\tilde{\mathbf{Z}}$ are produced by
row-shuffling $\mathbf{X}$ before encoding. The DGI loss is:
\begin{equation}
  \mathcal{L}_\text{DGI} = -\frac{1}{2N}
    \sum_{i=1}^{N}\!\bigl[
      \log\sigma\!\bigl(\mathcal{D}(\mathbf{z}_i, \mathbf{s})\bigr)
      + \log\!\bigl(1 - \sigma\!\bigl(
          \mathcal{D}(\tilde{\mathbf{z}}_i, \mathbf{s})\bigr)\bigr)
    \bigr].
  \label{eq:dgi}
\end{equation}
Pre-training runs for 300 epochs using Adam with learning
rate $10^{-3}$ and weight decay $5 \times 10^{-4}$.
The encoder weights are frozen upon completion and the
pre-computed embedding matrix $\mathbf{Z}$ is used as a
fixed skip-connection input to the downstream GAT,
following the transfer learning paradigm validated for
graph-based fraud detection by~\citet{graphanomaly2024}.

\subsubsection{Phase 2: Graph Attention Network Fine-Tuning}

The MixTrace GAT processes the full 169-dimensional feature
matrix through the following sequence:
\begin{align}
  \mathbf{H}_0 &= \text{Linear}(\mathbf{X})
    \in \mathbb{R}^{N \times 128}, \\
  \mathbf{H}_1 &= \text{BN}\!\bigl(\text{ELU}\!\bigl(
    \text{GAT}^{(K=8)}(\mathbf{H}_0, \mathcal{E})\bigr)\bigr)
    \in \mathbb{R}^{N \times 1024}, \\
  \mathbf{H}_2 &= \text{BN}\!\bigl(\text{ELU}\!\bigl(
    \text{GAT}^{(K=1)}(\mathbf{H}_1, \mathcal{E})\bigr)\bigr)
    \in \mathbb{R}^{N \times 128},
\end{align}
where $\text{GAT}^{(K)}$ denotes a $K$-head attention layer
and BN denotes batch normalisation.
Dropout with probability 0.3 follows each activation.
The frozen DGI embedding is concatenated to form a
256-dimensional fused representation:
\begin{equation}
  \mathbf{H}_\text{fused} = [\mathbf{H}_2 \;\|\; \mathbf{Z}]
    \in \mathbb{R}^{N \times 256}.
\end{equation}
A two-layer MLP classifier maps fused representations to
logits $\hat{\mathbf{y}} \in \mathbb{R}^{N \times 2}$.

\subsubsection{Training Protocol}

Fine-tuning uses class-weighted cross-entropy:
\begin{equation}
  \mathcal{L}_\text{GAT}
    = -\sum_{v \in \mathcal{V}_L}
      w_{y_v} \log \hat{p}(y_v \mid v),
\end{equation}
where $\mathcal{V}_L$ is the labelled training set and
$w_y$ is the inverse class-frequency weight.
Given 4,545 illicit and 42,019 licit training labels, the
weights are $w_0 \approx 0.565$ (licit) and
$w_1 \approx 4.317$ (illicit). The optimiser is Adam with
initial learning rate $10^{-3}$ and weight decay
$5 \times 10^{-4}$. A ReduceLROnPlateau scheduler monitors
validation F1-score with patience 10 and reduction factor 0.5.
Training terminates after 500 epochs or 20 consecutive
epochs without validation improvement.
ADASYN oversampling is applied to the labelled node
feature matrix as a supplementary diagnostic, following the
density-sensitive synthesis protocol of~\citet{tekkali2024adasyn}.

% ─────────────────────────────────────────────────────────────
\section{Experiments}
\label{sec:exp}
% ─────────────────────────────────────────────────────────────

\subsection{Datasets}

\paragraph{Elliptic Bitcoin Dataset.}
The Elliptic dataset comprises 203,769 Bitcoin transaction
nodes across 49 weekly time steps with 165 structural and
aggregated features. The class distribution among 46,564
labelled nodes is 4,545 illicit (9.76 percent) and 42,019
licit (90.24 percent), yielding a 9.2:1 imbalance ratio.
The directed edge list contains 234,355 entries. A strict
chronological split assigns time steps 1 to 34 to training,
35 to 41 to validation, and 42 to 49 to the held-out test
set, preserving temporal ordering to prevent data leakage.
The test partition contains 408 illicit and 8,433 licit
labelled nodes.

\paragraph{Author-Curated Bitcoin Transaction Corpus.}
A corpus of 5,884,387 Bitcoin transactions was assembled with
53 features including script-type indicators, fee rates, input
and output counts, address-reuse flags, value metrics, and
temporal features. The ground-truth label \texttt{is\_coinjoin\_like},
derived from structural heuristics validated against known
mixing-service transaction patterns, identifies 110,352
transactions (1.88 percent) as probable CoinJoin participants.
The dataset predates widespread P2TR adoption: no Taproot
transactions are present, and multi-input P2WPKH transactions
serve as the structurally equivalent CIOH proxy for Gap 1
analysis, consistent with the fallback strategy described
in Section~\ref{sec:taproot_method}.

\subsection{Baseline Methods}

\textbf{K-Means} ($k=2$) applies Euclidean clustering to the
standardised node feature matrix. Cluster-to-class assignment
uses majority vote over training labels.

\textbf{K-Nearest Neighbours (KNN)} applies grid search over
$k \in \{3,5,7,11,15\}$ and uniform/distance weighting,
selecting the configuration maximising validation illicit F1.

\textbf{Random Forest (RF)} applies grid search over estimator
count in $\{100,200\}$, maximum depth in $\{10,20,\text{None}\}$,
and minimum split samples in $\{2,5\}$, with balanced class
weighting throughout.

\textbf{GAT without DGI} reproduces the full MixTrace GAT
architecture with the DGI skip connection replaced by a zero
tensor, isolating the pre-training contribution from the
attention mechanism, consistent with ablation methodology
recommended by~\citet{motie2024gnn}.

\subsection{Evaluation Metrics}

The primary metric is the illicit-class F1-score.
Secondary metrics are precision, recall, macro F1, PR-AUC,
and ROC-AUC. For the Taproot analysis, false-positive rate,
false-negative rate, and DBSCAN silhouette score are reported.
For CST linking, the optimal threshold and corresponding F1,
precision, and recall are reported for both formulations.

\subsection{Illicit Transaction Detection Results}
\label{sec:classification_results}

Table~\ref{tab:classification} reports test-set performance
on the Elliptic held-out partition.

\begin{table}[htbp]
  \caption{Illicit transaction detection results on the
    Elliptic test set (time steps 42--49, 408 illicit and
    8,433 licit labelled nodes). Best result per column
    in \textbf{bold}.}
  \label{tab:classification}
  \centering
  \small
  \setlength{\tabcolsep}{5pt}
  \begin{tabular}{lcccccc}
    \toprule
    \textbf{Method} &
    \textbf{Prec.} & \textbf{Rec.} &
    \textbf{F1} & \textbf{F1\textsubscript{mac}} &
    \textbf{PR-AUC} & \textbf{ROC-AUC} \\
    \midrule
    K-Means
      & 0.046 & 1.000 & 0.088 & 0.045 & 0.000 & -- \\
    KNN
      & 0.499 & 0.468 & 0.483 & 0.729 & 0.416 & -- \\
    Random Forest
      & \textbf{0.946} & 0.471 & 0.628 & 0.808 & 0.632 & -- \\
    GAT (no DGI)
      & 0.113 & \textbf{0.995} & 0.204 & 0.486 & 0.269 & 0.932 \\
    \textbf{MixTrace}
      & 0.709 & 0.760 & \textbf{0.734} & \textbf{0.860} &
        \textbf{0.717} & \textbf{0.987} \\
    \bottomrule
  \end{tabular}
\end{table}

MixTrace achieves the highest illicit F1-score (0.7337),
macro F1 (0.8602), PR-AUC (0.7167), and ROC-AUC (0.9868).
Random Forest attains the highest precision (0.9458) at
the cost of recall (0.4706), producing a substantially
lower F1 of 0.6285. The standalone GAT without DGI exhibits
a degenerate high-recall, low-precision pattern (recall 0.995,
precision 0.113, F1 0.204), confirming that graph attention
layers alone, without structural pre-training context,
collapse to near-unconditional positive prediction under
the 9.2:1 imbalance ratio. K-Means exhibits the same
degenerate pattern at even lower precision, validating
that unsupervised clustering without neighbourhood information
is insufficient for minority-class detection at this
imbalance severity~\citep{albalawi2025credit}.

\subsection{DGI Pre-Training and Fine-Tuning Dynamics}

The DGI encoder was trained for 300 epochs. The initial
BCE loss at epoch 1 was 1.459, converging to stable low
values by epoch 300. At GAT fine-tuning epoch 5, training
F1 reached 0.8423 and validation F1 reached 0.5094. By
epoch 20, training F1 was 0.8844 and validation F1 improved
to 0.5847. Convergence stabilised in the range of epochs
220 to 265, where training F1 exceeded 0.9550 and the best
recorded validation macro F1 of 0.8864 was observed at epoch
235. The optimal decision threshold selected on the validation
set was $\theta^* = 0.9397$, indicating well-calibrated
high-confidence predictions rather than low-probability
borderline scores.

\subsection{Taproot Analysis Results}

Table~\ref{tab:taproot} presents the comparison between
naive CIOH and Taproot-Aware DBSCAN on the 200,000-transaction
analysis set. The set contains 22,169 CoinJoin-like and
177,831 non-CoinJoin-like transactions.

\begin{table}[htbp]
  \caption{Taproot-proxy analysis on 200,000 multi-input
    Bitcoin transactions. Sil.: DBSCAN silhouette score.}
  \label{tab:taproot}
  \centering
  \small
  \setlength{\tabcolsep}{4pt}
  \begin{tabular}{lccccccccc}
    \toprule
    \textbf{Method}
      & \textbf{TP} & \textbf{TN}
      & \textbf{FP} & \textbf{FN}
      & \textbf{FPR} & \textbf{FNR}
      & \textbf{Prec.} & \textbf{F1} & \textbf{Sil.} \\
    \midrule
    Naive CIOH
      & 22,169 & 171,208 & 6,623 & 0
      & 0.0372 & 0.000 & 0.770 & 0.870 & -- \\
    \textbf{DBSCAN}
      & 6,650 & 176,494 & 1,337 & 15,519
      & \textbf{0.0075} & 0.700 & \textbf{0.833} & 0.441
      & \textbf{0.973} \\
    \bottomrule
  \end{tabular}
\end{table}

Taproot-Aware DBSCAN reduces the false-positive rate from
3.72 to 0.75 percent: an absolute reduction of 0.0297 and
a relative reduction of 79.8 percent, eliminating over
5,200 incorrect attributions across the analysis set.
Precision improves from 0.77 to 0.8326. The silhouette
score of 0.973 confirms highly compact, well-separated
clusters in the six-dimensional structural feature space,
validating the hypothesis that CoinJoin and non-CoinJoin
multi-input transactions occupy geometrically distinct
regions of the feature landscape~\citep{chen2025multidist}.
The DBSCAN assignment produces a higher false-negative rate
(70 percent) because structurally atypical CoinJoin transactions
lie outside the dense cluster cores at the chosen $\varepsilon$
radius. This precision-recall trade-off is consistent with
forensic deployment contexts in which false accusations
carry higher legal and reputational cost than missed
detections~\citep{futureinternet2025}.

\subsection{CST Linking via Chamfer Distance}

Table~\ref{tab:chamfer} compares the 1-D and MD Chamfer
Distance formulations across 2,247 temporal windows and
5,000 evaluated window pairs.

\begin{table}[htbp]
  \caption{CoinJoin Spending Transaction linking performance
    at the optimal threshold $\theta^*$ maximising F1-score
    on the validation partition.}
  \label{tab:chamfer}
  \centering
  \small
  \begin{tabular}{lccccc}
    \toprule
    \textbf{Method} & \textbf{Dims}
      & $\boldsymbol{\theta^*}$
      & \textbf{F1} & \textbf{Prec.} & \textbf{Rec.} \\
    \midrule
    1-D Chamfer (timestamp)
      & 1 & 2.008 & 0.251 & 0.145 & 0.956 \\
    MD-Chamfer ($d=4$)
      & 4 & 3.384 & 0.245 & 0.141 & 0.954 \\
    \bottomrule
  \end{tabular}
\end{table}

Both formulations achieve recall above 0.954, confirming that
Chamfer-based linking captures the majority of true
CoinJoin-spend pairs. Precision is constrained by the large
number of structurally similar transaction windows in the
dataset (0.145 and 0.141, respectively). The absolute
F1-score difference between formulations is 0.006, within
the range of threshold-selection variability. This result
establishes an important boundary condition for MD-CD:
when the temporal feature alone provides near-ceiling recall,
the marginal contribution of structural dimensions is
attenuated because fee rate and input-output counts are
partially correlated with timestamp in coordinated mixing
activity. The MD-CD nonetheless operates at a substantially
higher optimal threshold (3.384 versus 2.008), reflecting
that the four-dimensional distance metric spans a broader
geometric scale and remains sensitive to structural
mismatches between non-contemporaneous window pairs,
a property that becomes more discriminative when post-mix
spending exhibits greater temporal dispersion.

% ─────────────────────────────────────────────────────────────
\section{Discussion}
\label{sec:discuss}
% ─────────────────────────────────────────────────────────────

\subsection{Contribution of DGI Pre-Training}

Isolating the DGI skip connection through the ablation
experiment reveals a 36-percentage-point F1 improvement
from 0.204 to 0.734. The standalone GAT learns discriminative
neighbourhood representations as evidenced by its ROC-AUC
of 0.932, but without the global structural context from
pre-training, the classifier defaults to near-unconditional
positive prediction, producing recall approaching unity at
precision below 0.12. The DGI embeddings pre-trained on all
203,769 nodes encode global topological context that local
two-hop message passing cannot recover from labelled nodes
alone. This finding aligns with the self-supervised graph
learning analysis of~\citet{graphanomaly2024}, who identified
that prototype representations learned from unlabelled
transaction nodes improve classifier calibration by up to
40 percentage points on severely imbalanced benchmarks.

The pre-training strategy effectively leverages the 157,205
unlabelled Elliptic nodes: structural regularities encoded in
the unlabelled majority translate into better-calibrated
illicit predictions once fine-tuned on the labelled minority.
This is a practically significant finding for the blockchain
forensics domain, where confirmed illicit labels are scarce
relative to the total transaction volume and where acquiring
additional labels requires costly investigative resources.

\subsection{Comparative Analysis with Baseline Classifiers}

Random Forest achieves the highest precision among all
evaluated methods (0.946) but recall of only 0.471, identifying
a high-confidence subset of illicit transactions while missing
more than half the true illicit test nodes. For financial
intelligence and compliance applications, the recall shortfall
of Random Forest represents unmonitored fund flows. MixTrace
achieves competitive precision (0.709) alongside substantially
higher recall (0.760), producing an F1 improvement of 10.5
percentage points over Random Forest and 25.1 points over KNN.
The PR-AUC improvement from 0.632 to 0.717 indicates superior
ranking quality of MixTrace's illicit-class scores across the
full precision-recall curve, consistent with the comparative
GNN benchmarking conducted by~\citet{fan2024gnn}.

The KNN baseline achieves nearly balanced precision and
recall (0.499 and 0.468) but at substantially lower values
than MixTrace, reflecting that flat feature-space distances
between transaction feature vectors do not capture the
structured neighbourhood patterns that graph-based methods
exploit. K-Means exhibits a degenerate all-positive prediction
bias with precision of only 0.046, confirming that
unsupervised clustering without graph structure cannot
distinguish illicit transaction clusters at this imbalance
level~\citep{xiao2025biEGCN}.

\subsection{False-Positive Reduction in Taproot Clustering}

The 79.8 percent relative reduction in false-positive rate
achieved by Taproot-Aware DBSCAN addresses a forensically
critical failure mode. Each CIOH false positive represents
a legitimate transaction incorrectly attributed to a mixing
entity. In law-enforcement and exchange compliance contexts,
false attributions cause account freezes, regulatory reporting
obligations, and misallocation of investigative resources.
The structural clustering approach eliminates over 5,200
incorrect attributions across the 200,000-transaction sample
compared to naive CIOH, at the cost of additionally missing
15,519 CoinJoin transactions that lie outside dense cluster
cores. The silhouette score of 0.973 validates that the
six-dimensional structural feature space creates genuinely
separable cluster geometry rather than statistical
artefacts~\citep{lee2024censor}. As P2TR adoption continues
to grow following Taproot activation, the proportion of
false positives under naive CIOH will increase; the
Taproot-Aware DBSCAN provides a structurally principled
remedy that scales with transaction volume through the
subsample-and-propagate architecture.

\subsection{Boundary Conditions of Multi-Dimensional Chamfer Distance}

The 0.006 absolute F1 difference between 1-D and MD-CD
requires careful methodological interpretation.
Post-mix spending in automated wallets clusters tightly in
time, creating a near-ceiling recall signal from timestamps
alone at the dataset's temporal resolution. The structural
dimensions in the MD-CD nonetheless identify a distinct
optimal threshold (3.384 versus 2.008) and contribute
complementary geometric information that is expected to
become more discriminative in datasets with higher within-class
temporal variance, for instance when time-locked scripts
or multi-round mixing protocols disperse post-mix spending
across multiple hours or days. The MD-CD formulation is
mathematically correct and theoretically principled; the
limited differential in the present experimental setting
reflects the dominance of the temporal prior in this
specific dataset rather than a general limitation of the
metric design~\citep{chen2025multidist}.

\subsection{Novel Feature Discriminative Analysis}

The shortest-path-to-illicit feature achieves a Pearson
correlation of $r = -0.9182$ with the ground-truth label,
substantially stronger than any feature in the original
165-dimensional Elliptic representation. This strong
negative correlation confirms that illicit transactions
form tight interconnected subgraphs, placing neighbouring
nodes within zero or one hop of confirmed illicit sources.
Licit transactions appear at the full BFS cap of six hops,
creating a near-binary distributional separation.
The illicit-neighbour ratio achieves a complementary
correlation of $r = +0.4159$, carrying independent predictive
signal beyond path distance. Both features provide
discriminative information entirely absent from the
original Elliptic feature set, and computing them requires
only the public edge list with no additional data collection,
making the augmentation immediately reproducible by other
researchers.

\subsection{Implications for Blockchain Forensics Policy}

The MixTrace framework demonstrates that combining
self-supervised graph representation learning with
multi-head attention classification substantially improves
illicit transaction detection over conventional supervised
baselines without requiring additional labelled data.
This is directly relevant to the Financial Action Task Force
guidance on virtual asset service providers, which requires
exchanges to implement risk-based transaction monitoring
systems capable of identifying mixer usage and suspicious
spending patterns~\citep{lo2023inspection,ahmed2024federated}.
The Taproot-Aware DBSCAN addresses a statutory compliance
gap: as P2TR adoption grows, CIOH-based automatic monitoring
systems will generate increasing false-positive alerts
that consume compliance officer capacity and potentially
trigger unjustified customer actions.
The federated learning extension evaluated by~\citet{baabdullah2024federated}
suggests a deployment architecture in which multiple
exchanges jointly train MixTrace's graph encoder while
keeping transaction data on-premises, enabling industry-wide
graph coverage without centralising sensitive customer data.

\subsection{Limitations}

Several limitations should be acknowledged. The author-curated
dataset contains no confirmed Taproot transactions, requiring
multi-input transactions as a structural proxy for the Gap~1
analysis. DGI pre-training requires the full transaction graph
in memory simultaneously, currently limiting scalability to
approximately 500,000 nodes on 64~GB RAM without mini-batch
extensions. The MD-CD computation is $O(mn)$ per window pair,
which may become prohibitive at exchange throughput scale
without approximate nearest-neighbour optimisation.
The analysis covers only on-chain Bitcoin data; off-chain
activity through the Lightning Network and cross-chain bridges
falls outside the scope of the current pipeline.

\subsection{Ethical Considerations}

The entire analysis is conducted on publicly available
pseudonymous on-chain data. No personally identifiable
information is processed or stored at any stage of the
pipeline. The Elliptic dataset labels were provided by a
regulated blockchain analytics provider whose methodology
has been independently reviewed~\citep{lo2023inspection}.
The intended application of the methodology is to support
financial intelligence, exchange compliance, and law-enforcement
practitioners in identifying potential illicit flows,
consistent with anti-money laundering regulatory frameworks
applicable to virtual asset service providers.

% ─────────────────────────────────────────────────────────────
\section{Conclusion}
\label{sec:conclude}
% ─────────────────────────────────────────────────────────────

This study presents MixTrace, a multi-modal graph attention
framework for Bitcoin CoinJoin transaction deanonymization
that simultaneously addresses three research gaps identified
in the current literature.

The Taproot-Aware DBSCAN strategy achieves a 79.8 percent
relative reduction in false-positive clustering rate over
naive CIOH across 200,000 analysed transactions, with a
silhouette score of 0.973 confirming high cluster quality.
The Multi-Dimensional Chamfer Distance extends the CST
linking metric to four structural dimensions; experimental
results establish the boundary conditions under which this
extension yields measurable F1 improvement over the
temporal baseline. The DGI plus GAT architecture achieves
an illicit-class F1-score of 0.7337, ROC-AUC of 0.9868,
and a 10.5 percentage-point F1 improvement over the strongest
conventional baseline. The ablation study demonstrates that
the DGI skip connection contributes 36 percentage points of
the final F1, confirming that self-supervised pre-training
on unlabelled blockchain nodes is the primary performance driver.

Three directions for future research are identified.
First, inductive DGI variants supporting mini-batch processing
would extend the framework to streaming blockchain data at
exchange scale~\citep{chithanuru2025ethereum}.
Second, integration with Lightning Network routing graphs
would enable off-chain mixing detection alongside on-chain
analysis. Third, federated graph training across multiple
exchange compliance systems would scale MixTrace's coverage
without centralising sensitive transaction data, following
the architecture proposed by~\citet{ahmed2024federated}.

The complete 14-step reproducible implementation pipeline
is available at \url{https://github.com/sagarkorde/CoinJoin}
under the MIT License.

% ─────────────────────────────────────────────────────────────
\section*{Acknowledgements}

The author thanks the Elliptic team for releasing the Bitcoin
benchmark dataset, the PyTorch Geometric and scikit-learn
open-source communities, and the guest editors of the special
issue on AI Enabled Security Frameworks for Emerging Cybercrime
and Digital Threats for the invitation to contribute.

\section*{Declarations}

\noindent\textbf{Conflict of interest.} The author declares
no conflict of interest.

\noindent\textbf{Funding.} No external funding was received.

\noindent\textbf{Data availability.} The Elliptic Bitcoin
Dataset is publicly available at
\url{https://www.kaggle.com/datasets/ellipticco/elliptic-data-set}.
The author-curated dataset is available on reasonable request.

\noindent\textbf{Code availability.} Full pipeline available at
\url{https://github.com/sagarkorde/CoinJoin} (MIT License).

% ─────────────────────────────────────────────────────────────
\bibliography{references}
\end{document}
